\documentclass[11pt,a4paper]{article}

\usepackage[margin=1in]{geometry}
\usepackage{amsmath}
\usepackage{microtype}
\usepackage{enumitem}
\usepackage{booktabs}
\usepackage{xcolor}
\usepackage{listings}
\usepackage{hyperref}

\hypersetup{
  colorlinks=true,
  linkcolor=black,
  urlcolor=blue!60!black,
  citecolor=black,
}

\definecolor{codebg}{HTML}{F6F8FA}
\definecolor{codecomment}{HTML}{6A737D}
\definecolor{codekw}{HTML}{D73A49}
\definecolor{codestr}{HTML}{032F62}

\lstset{
  basicstyle=\ttfamily\small,
  backgroundcolor=\color{codebg},
  commentstyle=\color{codecomment}\itshape,
  keywordstyle=\color{codekw}\bfseries,
  stringstyle=\color{codestr},
  showstringspaces=false,
  breaklines=true,
  frame=single,
  framesep=4pt,
  rulecolor=\color{codebg},
  language=Python,
  numbers=none,
  xleftmargin=4pt,
  xrightmargin=4pt,
}

\title{\textbf{CSL GraphBuilder --- Architectural Walk-through}\\
  \large Coroutines, Threads, Queues, and Pools in the Backend}
\author{}
\date{}

\begin{document}
\maketitle

\section{Overview}

CSL GraphBuilder is a biomedical knowledge-graph construction service. It
ingests documents (or external sources such as Open~Targets), extracts
entities and relationships, embeds them with a domain-specific
sentence-transformer (SapBERT), persists everything to Neo4j, and then runs
a cascading verification pipeline before exposing the graph to a Next.js
frontend.

The system is split into three deployable units:

\begin{itemize}[leftmargin=*]
  \item \textbf{Frontend} --- Next.js dev server on port \texttt{3000}
        (\texttt{frontend/}).
  \item \textbf{Backend} --- FastAPI/Uvicorn on port \texttt{8000}
        (\texttt{api/}).
  \item \textbf{Database} --- Neo4j with native vector indexes for
        embedding-based retrieval.
\end{itemize}

The startup script that wires the first two together is
\texttt{start.sh}; it cleans up listeners on ports \texttt{3000}/\texttt{8000}
(filtered with \texttt{lsof -sTCP:LISTEN} so only servers, not browser
client connections, are killed), activates the Python venv, and launches
both processes.

\section{Backend Layered Architecture}

The Python backend follows a hexagonal layering. Within
\texttt{src/graphbuilder/}:

\begin{itemize}[leftmargin=*]
  \item \texttt{domain/models/} --- pure dataclasses
        (\texttt{KnowledgeGraph}, \texttt{GraphRelationship},
        \texttt{DocumentChunk}). No I/O, no async.
  \item \texttt{core/} --- domain services that don't own state:
        \texttt{processing/semantic\_chunker.py},
        \texttt{verification/cascading.py},
        \texttt{verification/embedding.py},
        \texttt{verification/llm\_verifier.py}.
  \item \texttt{application/use\_cases/} --- orchestrators that compose
        domain services with infrastructure ports
        (\texttt{document\_pipeline.py},
        \texttt{relationship\_verification.py},
        \texttt{open\_targets\_ingestion.py}, ...).
  \item \texttt{infrastructure/} --- adapters: Neo4j repository
        (\texttt{repositories/graph\_repository.py}), embedding factory
        (\texttt{services/embedding\_factory.py}), GPU pool
        (\texttt{services/gpu\_embedding\_pool.py}), LLM service
        (\texttt{services/llm\_service.py}), config
        (\texttt{config/settings.py}).
  \item \texttt{cli/main.py} --- Click CLI for offline batch processing.
\end{itemize}

The HTTP surface lives in \texttt{api/main.py} and the routers under
\texttt{api/routers/} (\texttt{ingest.py}, \texttt{graph.py},
\texttt{verification.py}, \texttt{documents.py}, \texttt{curation.py},
\texttt{health.py}, \texttt{export.py}, \texttt{dev.py}). Routers are
thin: they parse the request, call an application use-case, and shape
the response.

\section{The Three Concurrency Primitives}

The backend deliberately mixes three forms of concurrency, each chosen
for the workload it is best at. Understanding which primitive owns which
job is the key to reading the codebase.

\subsection{Coroutines (asyncio) --- request and pipeline orchestration}

Every router handler is \texttt{async}. The application layer is
\texttt{async} all the way down: use-cases stage their work as
\texttt{await}-able steps, with a \texttt{cancel\_check} barrier between
each stage so that long-running pipelines can be aborted cleanly when a
user cancels a job. The canonical example is
\texttt{DocumentPipelineUseCase} in
\texttt{src/graphbuilder/application/use\_cases/document\_pipeline.py},
which executes:

\begin{enumerate}[leftmargin=*,itemsep=2pt]
  \item \texttt{\_stage\_chunk}     --- semantic chunking
  \item \texttt{\_stage\_entities}  --- entity extraction
  \item \texttt{\_stage\_relationships} --- relationship extraction
  \item \texttt{\_stage\_persist}   --- save to Neo4j
  \item \texttt{\_stage\_verify}    --- cascading verification
\end{enumerate}

Coroutines own \emph{orchestration}, not CPU-bound work. They are cheap
to create, they preserve a single shared event loop, and they let the
pipeline issue many concurrent network and database calls (e.g.\ Neo4j
batch writes, LLM calls) without spawning OS threads.

\paragraph{Why this matters for the LLM call.}
\texttt{LLMVerifier.verify} in
\texttt{src/graphbuilder/core/verification/llm\_verifier.py} is
\texttt{async} and awaits the underlying \texttt{llm\_service.generate\_text}.
Many such verifications can be in flight on the same loop at once,
limited only by the \texttt{VERIFY\_PARALLEL\_WORKERS} setting in
\texttt{src/graphbuilder/infrastructure/config/settings.py}.

\subsection{Threads --- where blocking math happens}

The event loop must \emph{never} block. Two situations force us into
threads:

\begin{enumerate}[leftmargin=*]
  \item \textbf{PyTorch / sentence-transformers inference.}
        \texttt{SentenceTransformer.encode} is a synchronous CUDA-or-CPU
        call that takes tens to hundreds of milliseconds. Calling it
        directly from an async handler freezes \texttt{/health} and
        every other in-flight request.
  \item \textbf{Model warm-up and pool initialisation.}
        SapBERT is $\sim$440~MB and takes 5--15~s to load. Doing this
        synchronously at FastAPI startup would delay the first request.
\end{enumerate}

The threading lives in
\texttt{src/graphbuilder/infrastructure/services/gpu\_embedding\_pool.py}
and \texttt{src/graphbuilder/infrastructure/services/embedding\_factory.py}:

\begin{itemize}[leftmargin=*]
  \item One \emph{init thread} (\texttt{daemon=True},
        \texttt{name="GPUEmbeddingPool-init"}) loads + warms a base
        model, measures per-worker GPU memory, then spawns the worker
        threads. The constructor returns immediately so FastAPI startup
        is not blocked.
  \item $N$ \emph{worker threads}
        (\texttt{name="GPUEmbeddingWorker-i"}), each owning its own
        \texttt{SentenceTransformer} replica and a dedicated
        \texttt{torch.cuda.Stream} so the GPU scheduler can interleave
        kernels instead of serialising on the default stream.
  \item On \emph{CPU-only} machines the pool is never instantiated; a
        single \texttt{asyncio.Lock} plus the default
        \texttt{run\_in\_executor} thread are used (see
        \texttt{embed\_async} / \texttt{embed\_batch\_async} in
        \texttt{embedding\_factory.py}). The lock is required because
        PyTorch CPU inference on the same model is not safe under
        concurrent access; on CPU there is also nothing to gain from
        parallel encodes because MKL/OpenMP already saturates cores
        from a single call.
\end{itemize}

\subsection{Queue --- the GPU work-stealing pool}

The single coupling between the asyncio world and the worker threads is
a \texttt{queue.Queue} of \texttt{\_Job} items inside
\texttt{GPUEmbeddingPool}:

\begin{lstlisting}[language=Python]
@dataclass
class _Job:
    payload: Union[str, List[str]]
    is_batch: bool
    batch_size: int
    future: asyncio.Future
    loop:   asyncio.AbstractEventLoop
\end{lstlisting}

The async API (\texttt{embed\_async}, \texttt{embed\_batch\_async}) creates
an \texttt{asyncio.Future} on the calling loop, packages it into a
\texttt{\_Job}, drops it on the queue, and \texttt{await}s the future.
Worker threads pull from the queue with a 0.5~s timeout (so they can
notice the \texttt{stop\_event} during shutdown), run the encode under
their CUDA stream, and resolve the future with
\texttt{loop.call\_soon\_threadsafe}---the result lands back on the
event-loop thread cleanly, with no asyncio primitives touched
off-loop.

The load-balancing policy is \emph{whoever is free first grabs the
next job}. This is the simplest correct algorithm and matches what the
GPU actually wants: a busy worker does not get a new job, an idle one
does. Per-worker counters (\texttt{jobs\_in\_flight}) are tracked under
a \texttt{threading.Lock} for observability.

\subsection{Pools --- one model replica per worker}

The pool's worker count is sized once at start-up. After loading and
warming a single model, peak GPU memory is read with
\texttt{torch.cuda.max\_memory\_allocated} (this captures weights
\emph{plus} activations for a 32-batch forward pass), then:

\[
n = \min\!\left(\, \mathrm{max\_workers},\;
        \max\!\left(\mathrm{min\_workers},\;
        \left\lfloor
          \tfrac{\mathrm{free\_bytes}\;\cdot\;\mathrm{memory\_fraction}}
                {\mathrm{per\_worker\_bytes}}
        \right\rfloor
        \right) \right)
\]

with defaults \texttt{min\_workers=1}, \texttt{max\_workers=8},
\texttt{memory\_fraction=0.7}. Each of the resulting $n$ workers loads
its own \texttt{SentenceTransformer} replica into VRAM. The number is
\textbf{static for the lifetime of the process}: there is no
auto-scaler, and \texttt{shutdown()} only signals workers to stop, it
does not resize.

Override knobs (read by \texttt{embedding\_factory.py}):

\begin{itemize}[leftmargin=*]
  \item \texttt{EMBEDDING\_GPU\_WORKERS} --- exact count, bypasses the
        formula above.
  \item \texttt{EMBEDDING\_GPU\_MIN\_WORKERS} / \texttt{\_MAX\_WORKERS} ---
        clamp range.
  \item \texttt{EMBEDDING\_GPU\_MEMORY\_FRACTION} --- fraction of free
        VRAM the pool may consume.
  \item \texttt{EMBEDDING\_GPU\_DEVICE} --- target device, default
        \texttt{cuda:0}.
\end{itemize}

\section{End-to-End Flow of an Embedding}

The following walk-through traces a single embedding request from the
chunker down to the GPU and back, illustrating how all four primitives
collaborate.

\begin{enumerate}[leftmargin=*,itemsep=2pt]
  \item A document arrives via
        \texttt{POST /documents} (router: \texttt{api/routers/documents.py}).
        FastAPI's event loop schedules the handler as a coroutine.
  \item The handler delegates to
        \texttt{DocumentPipelineUseCase.execute}
        (\texttt{src/graphbuilder/application/use\_cases/document\_pipeline.py}).
        \texttt{\_stage\_chunk} constructs a \texttt{SemanticChunker}
        and calls \texttt{await chunker.chunk(content, document.id)}.
  \item Inside \texttt{SemanticChunker.\_embed}
        (\texttt{src/graphbuilder/core/processing/semantic\_chunker.py})
        the sentence list is passed to
        \texttt{embed\_batch\_async(sentences)}.
  \item \texttt{embed\_batch\_async} in
        \texttt{embedding\_factory.py} probes for a GPU pool. If
        present, it forwards to
        \texttt{GPUEmbeddingPool.embed\_batch\_async}; if not, it takes
        the per-loop \texttt{asyncio.Lock} and runs \texttt{embed\_batch}
        in the default thread executor.
  \item \emph{(GPU path.)} The pool builds a \texttt{\_Job} with the
        coroutine's future and \texttt{put}s it on the worker queue.
        The coroutine \texttt{await}s the future.
  \item A free worker thread pops the job, switches to its dedicated
        \texttt{torch.cuda.Stream}, and runs \texttt{model.encode}.
  \item The worker resolves the future via
        \texttt{loop.call\_soon\_threadsafe(future.set\_result,
        result)}. The event loop's I/O thread re-enters the awaiting
        coroutine on the next tick.
  \item The chunker turns the returned \texttt{List[Optional[List[float]]]}
        into an \texttt{(N, D)} \texttt{np.ndarray}, runs the
        sentence-grouping algorithm, and returns
        \texttt{DocumentChunk}s up the call stack.
\end{enumerate}

The same pattern recurs for entity / relationship embeddings during
\texttt{save\_entities\_batch} and \texttt{save\_relationships\_batch}
in \texttt{src/graphbuilder/infrastructure/repositories/graph\_repository.py}.

\section{Verification Cascade}

A second concurrency hot spot is verification, in
\texttt{src/graphbuilder/core/verification/cascading.py}. Each
relationship goes through up to three stages, in order of cost:

\begin{enumerate}[leftmargin=*]
  \item \textbf{Text-match} --- substring/lexical, instant.
  \item \textbf{Embedding} --- cosine similarity against the
        relationship's source chunk; uses the same
        \texttt{embed\_batch\_async} path.
  \item \textbf{LLM} --- structured-JSON prompt to
        \texttt{LLMVerifier} (\texttt{llm\_verifier.py}), gated by
        \texttt{VERIFY\_BATCH\_SKIP\_LLM} (default \texttt{true}, so
        Stage 3 is skipped during pipeline batches and runs only when
        the user manually triggers it from the
        \texttt{/verification} page).
\end{enumerate}

The verification use-case
(\texttt{application/use\_cases/relationship\_verification.py}) iterates
relationships in coroutine form and dispatches the three stages
sequentially per relationship; concurrency between relationships is
controlled by the bounded
\texttt{VERIFY\_PARALLEL\_WORKERS} semaphore in the pipeline orchestrator.

\section{Configuration Surface}

All concurrency-related knobs live in
\texttt{src/graphbuilder/infrastructure/config/settings.py}:

\medskip
\begin{tabular}{@{}lll@{}}
  \toprule
  Setting / env var & Default & Effect \\
  \midrule
  \texttt{PROCESSING\_MODE}         & \texttt{batch}  & Selects bulk vs single-doc paths \\
  \texttt{VERIFY\_ENABLED}          & \texttt{true}   & Master switch for the verify stage \\
  \texttt{VERIFY\_BATCH\_SKIP\_LLM} & \texttt{true}   & Skip LLM stage in pipeline batches \\
  \texttt{VERIFY\_PARALLEL\_WORKERS}& \texttt{4}      & In-pipeline relationship concurrency \\
  \texttt{EMBEDDING\_GPU\_WORKERS}  & ---             & Force exact pool size \\
  \texttt{EMBEDDING\_GPU\_MIN\_WORKERS} & \texttt{1}  & Pool sizing floor \\
  \texttt{EMBEDDING\_GPU\_MAX\_WORKERS} & \texttt{8}  & Pool sizing ceiling \\
  \texttt{EMBEDDING\_GPU\_MEMORY\_FRACTION} & \texttt{0.7} & VRAM share for pool \\
  \texttt{EMBEDDING\_GPU\_DEVICE}   & \texttt{cuda:0} & Target CUDA device \\
  \texttt{EMBEDDING\_MODEL}         & SapBERT         & Sentence-transformer model name \\
  \bottomrule
\end{tabular}

\section{File Index}

\begin{tabular}{@{}p{0.45\linewidth}p{0.50\linewidth}@{}}
  \toprule
  Path & Role \\
  \midrule
  \texttt{start.sh}                                        & Process supervisor, port cleanup \\
  \texttt{api/main.py}                                     & FastAPI app + lifespan hooks \\
  \texttt{api/routers/*.py}                                & HTTP routers (ingest, graph, verify, ...) \\
  \texttt{application/use\_cases/document\_pipeline.py}    & Async ingest pipeline \\
  \texttt{application/use\_cases/relationship\_verification.py} & Verification orchestrator \\
  \texttt{application/use\_cases/open\_targets\_ingestion.py} & External-source ingestion \\
  \texttt{core/processing/semantic\_chunker.py}            & Sentence-aware async chunker \\
  \texttt{core/verification/cascading.py}                  & Three-stage verifier \\
  \texttt{core/verification/llm\_verifier.py}              & LLM JSON verifier \\
  \texttt{infrastructure/services/embedding\_factory.py}   & Backend-selection async API \\
  \texttt{infrastructure/services/gpu\_embedding\_pool.py} & Threaded GPU pool + queue \\
  \texttt{infrastructure/services/llm\_service.py}         & LLM client \\
  \texttt{infrastructure/repositories/graph\_repository.py}& Neo4j persistence \\
  \texttt{infrastructure/config/settings.py}               & Env-driven configuration \\
  \bottomrule
\end{tabular}

\section{Take-aways}

\begin{itemize}[leftmargin=*]
  \item \textbf{Coroutines own orchestration.} Every router and use-case
        is \texttt{async}; the event loop is the single coordination
        point.
  \item \textbf{Threads own blocking math.} GPU encodes and model loads
        run on threads so the loop never blocks; an
        \texttt{asyncio.Future} bridges the two worlds.
  \item \textbf{The queue decouples them.} \texttt{queue.Queue} is the
        only object shared between the loop and the workers; the
        ``free worker grabs next job'' policy is implicit
        load-balancing.
  \item \textbf{The pool is statically sized.} Worker count is decided
        once at boot from a peak-memory measurement and a fraction of
        free VRAM, then frozen for the process lifetime.
  \item \textbf{CPU mode degrades gracefully.} Without CUDA the pool is
        skipped entirely; a single shared model behind an
        \texttt{asyncio.Lock} plus \texttt{run\_in\_executor} keeps the
        loop responsive even though encodes are serialised.
\end{itemize}

\end{document}
